\documentclass[../main.tex]{subfiles}
\graphicspath{{\subfix{../image/}}}
\begin{document}
\chapter{淹没、浸入与嵌入}\label{chapter:ch04-淹没浸入嵌入}

由于光滑映射的微分被认为是该映射在给定点附近的``最佳线性逼近'',故通过研究其微分的线性代数性质,我们可以获得关于该映射本身的大量信息。微分最本质的性质——事实上,差不多是唯一与基的选取无关即可定义的性质——就是它的秩(即其像的维数)。

本章将详细研究:光滑映射的几何性质如何能从其微分中被识别出来。那些微分能给出良好局部模型的映射,恰好就是微分具有常秩的映射。这类映射中有三类扮演着特殊角色:光滑淹没(其微分处处满射)、光滑浸入(其微分处处单射),以及光滑嵌入(同时还是到其像上之同胚的单射光滑浸入)。正如我们将在下一章看到的,光滑浸入与光滑嵌入是子流形理论的基本要素;而光滑淹没在光滑流形理论中所起的作用,与商映射在拓扑学中所起的作用密切类似。

驱动整个讨论的引擎是秩定理,它是反函数定理的一个推论。在本章第一节,我们证明秩定理及其若干重要推论。随后我们更深入地研究光滑嵌入与光滑淹没,并把该理论应用于一类特别有用的光滑淹没,即光滑覆叠映射。

\subfile{Chapter-04/section-01.tex}
\subfile{Chapter-04/section-02.tex}
\subfile{Chapter-04/section-03.tex}
\subfile{Chapter-04/section-04.tex}
\subfile{Chapter-04/section-05.tex}
\end{document}